% Zdroj: PDF priklady-jaro-2026.pdf, fyzická str. 2-3. Značka: společné okruhy. Zařazeno: I4.
\section{Semafor mezi procesy}
\szzotazka{jaro 2026}{2}{Semafor mezi procesy (FIFO buffer)}

\noindent\textbf{Zadání.}\par
Implementujte dvojici funkcí s~tímto rozhraním:

\begin{lstlisting}[language=C]
void send(byte ch);
byte receive();
\end{lstlisting}

Účelem těchto funkcí je realizace jednosměrného proudu bajtů mezi dvěma procesy. Součástí
proudu je FIFO-buffer s~kapacitou $N$ bajtů. Funkce \texttt{receive}, volaná procesem
příjemce, čeká, dokud v~bufferu není alespoň jeden platný bajt, poté tento bajt odebere
a~vrátí jej. Funkce \texttt{send}, volaná z~procesu odesilatele, čeká, dokud v~bufferu
není alespoň jedna volná pozice, poté do bufferu vloží parametr \texttt{ch}.

Vaším úkolem je implementovat funkce \texttt{send} a~\texttt{receive} pouze pomocí těchto součástí:
\begin{itemize}
\item Semafory (pracující přes hranice procesů)
\item Implementace bufferu s~kapacitou $N$, která poskytuje následující rozhraní:
\end{itemize}

\begin{lstlisting}[language=C]
void buffer_push(byte ch);
byte buffer_pop();
\end{lstlisting}

Funkce \texttt{buffer\_push} (volaná procesem odesilatele) a~funkce \texttt{buffer\_pop}
(volaná procesem příjemce) jsou atomické (tj. mohou být bezpečně volány současně)
a~umožňují předávání bajtů přes hranice procesů (např. pomocí sdílené paměti a~atomických
instrukcí). Tyto funkce však neposkytují synchronizaci (vždy se vracejí okamžitě). Funkce
\texttt{buffer\_push} nesmí být volána, pokud je buffer plný, a~funkce \texttt{buffer\_pop}
nesmí být volána, pokud je buffer prázdný.

Vaším úkolem není řešit vytvoření semaforů a~bufferu, předpokládejte, že oba procesy mohou
přistupovat k~těmto objektům, jako by existovaly uvnitř těchto procesů.

\begin{enumerate}
\item Napište implementaci funkcí \texttt{send} a~\texttt{receive} v~některém z~jazyků C,
  C++, C\# nebo Java. Implementace nesmí používat žádné statické nebo globální proměnné
  ani žádné další knihovny kromě semaforů a~výše popsaného bufferu. Implementace nesmí
  využívat aktivní čekání (polling). Specifikujte také počáteční stav(y) semaforu/ů
  (v~závislosti na kapacitě bufferu $N$); buffer je na počátku prázdný.
\item Definujte invariant(y) platné pro stav(y) semaforu/ů ve vztahu k~počtu platných
  elementů v~bufferu (a~jeho kapacitě $N$). Všechny tyto invarianty musí být platné po
  celou dobu mezi voláními funkcí \texttt{send} nebo \texttt{receive}.
\item Vysvětlete, jak se platnost invariantů mění uvnitř funkcí \texttt{send}
  a~\texttt{receive} (např. pomocí pozměněných invariantů platných mezi jednotlivými
  příkazy uvnitř těchto funkcí). Pro jednoduchost předpokládejte, že tyto funkce nejsou
  volány zároveň.
\end{enumerate}

\medskip
\noindent\textbf{Náčrt řešení.}\par
\begin{lstlisting}[language=C]
// capacity_semaphore initialized to N
// content_semaphore initialized to 0
//
// Invariant 1: capacity_semaphore = N - elements_in_buffer
// Invariant 2: content_semaphore = elements_in_buffer

void send(char ch)
{
  capacity_semaphore.down();
  // Inv 1 skewed by +1, Inv 2 valid
  buffer_push(ch);
  // Inv 1 valid, Inv 2 skewed by -1
  content_semaphore.up();
}

char receive()
{
  content_semaphore.down();
  // Inv 1 valid, Inv 2 skewed by +1
  char ch = buffer_pop();
  // Inv 1 skewed by -1, Inv 2 valid
  capacity_semaphore.up();
}
\end{lstlisting}

\medskip
\noindent\textit{Doplňující vysvětlení (nad rámec oficiálního náčrtu):} V~náčrtu funkce
\texttt{receive} chybí na konci \texttt{return ch;}. Komentáře „skewed by $\pm1$`` říkají, o~kolik
se invariant mezi příkazy rozchází: po \texttt{capacity\_semaphore.down()} je hodnota semaforu o~1
menší, než by odpovídalo počtu prvků (jedno místo je už rezervované, ale ještě nezaplněné);
\texttt{buffer\_push} invariant~1 srovná a~invariant~2 je naopak o~1 pozadu, dokud ho
\texttt{content\_semaphore.up()} nesrovná. Funkce \texttt{receive} je symetrická. Další zámek není
potřeba, protože \texttt{buffer\_push} a~\texttt{buffer\_pop} jsou podle zadání atomické
a~semafory zaručují, že se push nevolá na plný a~pop na prázdný buffer.
